\documentclass{ctexart}
\usepackage{amsmath}
\usepackage{amssymb}
\usepackage{amsfonts}
\usepackage{xcolor}
\usepackage{graphicx}
\usepackage{booktabs}
\usepackage{array}
\usepackage{float}
\usepackage{hyperref}
\hypersetup{colorlinks=false}
\pagestyle{plain}
\ctexset{
    section = {name = {,、}, number = \chinese{section}},
    subsection = {name = {（,）}, number = \chinese{subsection}},
    subsubsection = {number = \arabic{subsubsection}}
}

\begin{document}


\section{测试转换}


\subsection{测试标题}

这是一个测试文档，用于验证 Markdown 转 LaTeX 的转换功能。

这是一个段落，包含一些\textbf{加粗}和\textit{斜体}文本，以及一个行内代码示例 \texttt{let x = 10;}。


\begin{verbatim}
fn main() {
    println!("Hello, world!");
}
\end{verbatim}

这是行内代码示例：\texttt{let x = 10;}，应该被正确转换为 LaTeX 的 \texttt{\textbackslash{}texttt\{let x = 10;\}}。

这是一个无序列表：


\begin{itemize}
\item 项目一
\item 项目二
\item 项目三
\end{itemize}

这是一个有序列表：


\begin{enumerate}
\item 第一项
\item 第二项
\item 第三项
\end{enumerate}


\subsection{数学公式测试}

这是行内公式：$E=mc^2$，应该原样保留。

这是块级公式：

\begin{equation*}
\int_{-\infty}^{\infty} e^{-x^2} dx = \sqrt{\pi}
\end{equation*}

这是第二个块级公式： \begin{equation*}
f(x) = \sum_{n=0}^{\infty} \frac{f^{(n)}(0)}{n!} x^n
\end{equation*}

给公式添加颜色： \begin{equation*}
\textcolor{red}{E=mc^2}
\end{equation*}

这是一个有编号的公式：


\begin{verbatim}
\begin{equation}
a^2 + b^2 = c^2
\end{equation}
\end{verbatim}

公式中的特殊字符如 \texttt{\textasciicircum{}}, \texttt{\_}, \texttt{\{}, \texttt{\}} 不应被转义。


\subsection{图片测试}


\begin{figure}[H]
\centering
\includegraphics[width=0.8\textwidth]{Gemini_Generated_Image_a40x9qa40x9qa40x.png}
\caption{项目印象图}
\end{figure}



下面是不同大小的图片测试：


\begin{figure}[H]
\centering
\includegraphics[width=0.8\textwidth]{Gemini_Generated_Image_a40x9qa40x9qa40x.png}
\caption{图片测试1}
\end{figure}


\begin{figure}[H]
\centering
\includegraphics[width=0.8\textwidth]{Gemini_Generated_Image_a40x9qa40x9qa40x.png}
\caption{图片测试2}
\end{figure}

其他的html width属性测试：


\begin{figure}[H]
\centering
\includegraphics[width=0.5\textwidth]{Gemini_Generated_Image_a40x9qa40x9qa40x.png}
\caption{图片测试3}
\end{figure}


\begin{figure}[H]
\centering
\includegraphics[width=0.8\textwidth]{Gemini_Generated_Image_a40x9qa40x9qa40x.png}
\caption{图片测试4}
\end{figure}


\subsection{链接测试}


\begin{itemize}
\item 这是一个链接：\href{https://www.openai.com}{点击这里访问 OpenAI}
\item 这是一个链接：\href{https://github.com}{点击这里访问 GitHub}
\item 这是一个链接：\href{https://www.rust-lang.org}{点击这里访问 Rust 官网}
\end{itemize}


\subsection{表格测试}

这是一个简单的表格：


\begin{table}[H]
\centering
\begin{tabular}{lcr}
\toprule
左对齐 & 居中 & 右对齐 \\
\midrule
内容1 & 内容2 & 内容3 \\
内容4 & 内容5 & 内容6 \\
\bottomrule
\end{tabular}
\end{table}

第二个表格：


\begin{table}[H]
\centering
\begin{tabular}{ll}
\toprule
姓名 & 年龄 \\
\midrule
张三 & 25 \\
李四 & 30 \\
\bottomrule
\end{tabular}
\end{table}

复杂表格：


\begin{table}[H]
\centering
\begin{tabular}{lccr}
\toprule
参数名称 & 符号 & 默认值 & 单位 \\
\midrule
采样频率 & $f_s$ & 1000 & Hz \\
截止频率 & $\omega_c$ & 50 & rad/s \\
阻尼比 & $\zeta$ & 0.707 & - \\
\bottomrule
\end{tabular}
\end{table}

带公式的表格：


\begin{table}[H]
\centering
\begin{tabular}{ll}
\toprule
参数 & 公式 \\
\midrule
频率响应 & $H(j\omega) = \frac{1}{1 + j\omega/\omega_c}$ \\
阻尼比 & $\zeta = \frac{c}{2\sqrt{mk}}$ \\
\bottomrule
\end{tabular}
\end{table}

带空行的表格：


\begin{table}[H]
\centering
\begin{tabular}{lll}
\toprule
实验编号 & 观测现象描述 & 备注 \\
\midrule
EXP-01 & 系统在输入阶跃信号后出现小幅震荡，超调量约 15\%。 & 正常 \\
EXP-02 &  & 数据待补 \\
EXP-03 & 调整 PID 参数后，稳态误差消除，收敛速度明显加快。 & \textbf{推荐参数} \\
\bottomrule
\end{tabular}
\end{table}



\end{document}
